\chapter{Objetivos} \label{sec:Objetivos}

\section{Objetivo General}

Desarrollar un modelo de evaluación, selección y validación de plugins CNI en Kubernetes, compuesto por una guía metodológica, un prototipo de pruebas de laboratorio y una lista de verificación.

\section{Objetivos Específicos}

\begin{enumerate}
  \item Evaluar cuantitativamente el rendimiento de los CNI disponibles para Kubernetes (Calico, Cilium, Flannel y Antrea) mediante pruebas controladas de latencia, velocidad de transferencia y uso de recursos, para identificar sus fortalezas y limitaciones en escenarios reales.
  \item Verificar la aplicación efectiva de políticas de red automatizadas en el plano de datos de cada plugin CNI y cuantificar el sobrecosto de rendimiento que introducen, con el fin de determinar qué plugins materializan realmente el aislamiento declarado y a qué costo, en entornos de laboratorio representativos de producción.
  \item Desarrollar, como parte del modelo propuesto, un conjunto de criterios y métricas que permitan la comparación objetiva entre diferentes CNI para Kubernetes mediante la definición de umbrales de rendimiento y factores de ponderación, para facilitar la toma de decisiones técnicas basadas en evidencia cuantificable.
  \item Implementar un prototipo funcional que recomiende y configure el CNI más adecuado según las necesidades específicas de cada organización, para simplificar el proceso de implementación y apoyar la definición de configuraciones recomendadas en entornos de prueba.
\end{enumerate}